\section{Electronic antecedent and dimensional realization}
\label{sec:electron}

The electronic construction receives the arithmetic state with
trajectory memory produced by APP--TRIT--TPK and its numerical coordinates. At this point in the development,
\(\pi,\varphi,e\) and \(\alpha\) designate the coordinates already obtained from
the same genealogy, in the constructive order set out above. They
are not parameters taken from a physical table. The electronic section is a
subsequent realization of that joint discrete structure of the continuum:
it preserves the cell of origin, the two arithmetic sheets, orientation,
carry, and return memory before producing a scalar.

Three objects should be distinguished from the outset. The first is the central
section and its transport history; the second is the dimensionless
coordinate obtained from it; the third is its realization on an
energy or mass line. Equality between two numerical coordinates does not
identify these three levels. In particular, Euler's number \(e\),
the electronic coordinate \(\varepsilon_e\), and a rest energy
\(E_e\) will have distinct notations.

The aim of this section is to establish the composition
\begin{equation}
 \varepsilon_e^{\beta}
 =\frac{\sqrt3}{4}
  \left[\exp\!\left(\frac{\varphi}{\pi^2}\right)-22\alpha^3\right]
  (1+15\Delta_4)
  R_{\mathrm{act}}^{\,80-54\alpha+6\Delta_4},
 \label{eq:el-formula-completa}
\end{equation}
explaining separately the origin of the prefactor, the two integers of the
basal stage, the normalization of \(\Delta_4\), the action ratio, and the
return triple. The dimensional chart is introduced only after
this composition. The result does not consist in making a
dimensionless number acquire units merely through its resemblance to a tabulated
quantity.

\paragraph{Identity of the section and composition of its observables.}
\label{el-ampl-articulacion}
The structural identity of the electron is introduced through the central
section and its transport operations. Inversion of the APP chart
selects the cell $(5,5)$ by its fixed-point property. The two arithmetic
sheets allow the deformations preserving the sum and the product
residue to be studied. The result distinguishes the two orientations
$(8,2)$ and $(2,8)$ and preserves a difference of one quotient unit.
This is the first level of information: a cell, an involution, a
preserved fibre, and a minimal integer defect. No mass evaluation
is required to state or prove it.

The passage to arithmetic projections introduces a second level. The
sum and product fibres determine conditional expectations on the
same finite space. Their commutator measures dependence on the
order of these two projections. The norm $\sqrt3/4$ follows from the spectrum
of their pairing; the block attaining this norm also preserves
the quadratic relations proved below. Thus, the scalar
prefactor and the local algebra have a common provenance, but
contain different information. The norm retains the magnitude of the
commutator, while the matrices preserve its orientation and its
composition law.

Areal--volumetric transfer incorporates the regional coordinates
at a third level. The mode calendar determines the incidence
matrix and its quadratic action. The closure mark produces the ratio
$\pi/\varphi^2$; interchanging the two arguments provides the
orientation $\varphi/\pi^2$ used in the electronic exponential
character. The two expressions are related by an explicit
involution. This relationship explains why a regional coordinate may
appear in the electronic equation with a power or in a position
different from that in its first reading: the operation applied
is part of the mathematical datum that is preserved.

The cubic alpha term belongs to the regional complement of a marked
return. Counting its positions yields $27-5=22$, while the
three sheets of each of the five regions yield $3\cdot5=15$.
The proofs preserve the domain of these counts and the
orientation determining their signs. The basal formula therefore
composes an incompatibility norm, a transfer character,
a position deficit, and a torsional memory.
Writing it on a single line expresses the composition of these
operations; explaining it requires reconstructing them in that order.

The transition from the basal coordinate to the complete coordinate preserves
the same central section. The return factor uses the history
records, and the two stated readers recover the triple
$(80,54,6)$ along that trajectory. Multiplication by
$R_{\mathrm{act}}^{80-54\alpha+6\Delta_4}$ modifies the evaluation,
but keeps its origin and antecedents identified. The basal
value constitutes an intermediate stage of the composition; the complete
value also includes this return contribution. The physical
comparison must use the stage specified by the formula being evaluated.

The action scale intervenes subsequently with a different function. A ratio of
two action sections is dimensionless and cancels their common scale.
The absolute realization instead uses an action basis, a
temporal reader, and the map to energy or mass. This separation allows
the construction of the scale and the identity of the electronic section
to be preserved simultaneously. The chronological quantum provides a
scale reference; the central section determines the structural factor
expressed relative to it. The dimensional equations of this
chapter keep that composition explicit.

The electron's relevance to the overall argument arises from this
convergence of operations. The same arithmetic that produced the
regions now determines a central cell, its deformations, and its
projections. The fine-structure constant, constructed earlier, acts
in this composition as a closure coordinate. Return information
determines the complete evaluation, and dimensional realization
allows subsequent comparison. Explaining the electron from its structure
consists in preserving this succession of objects, maps, and values at every
step of the proof.



\subsection{Center, orientation, and vacancy on a preserved fiber}

The cellular support under consideration is \(\{1,\ldots,9\}^2\), with both
APP evaluations, \(r+c\) and \(rc\). The central inversion is
\begin{equation}
 \rho(r,c)=(10-r,10-c).
 \label{eq:el-inversion-central}
\end{equation}
The center is not chosen by a mass: it is characterized by a property of
the support that precedes its physical reading.

\begin{proposition}[Uniqueness of the center]
The inversion \(\rho\) has a unique fixed point, \((5,5)\).
\end{proposition}
\begin{proof}
The equations \(r=10-r\) and \(c=10-c\) are equivalent to
\(2r=2c=10\). Their unique solution in the nonadic chart is \((5,5)\).
\end{proof}

The preceding property identifies a cell, not a unique deep
history. A route adds the initial choice of orientation and its
compatible prolongations. In the chart used here, the central
oriented section begins at \((5,5)\) with both orientations
positive; it will be denoted by \(\Gamma_e\). The return of its cellular
projection does not erase the transported orientations or the carry
record of the section with its trajectory memory.

Vacancy is now considered, after the closure of \(\alpha\), as
a local operation on this section. Its arithmetic can be isolated in a
finite lemma; this isolation does not make it a substitute for the
preceding generator of constants.

\begin{definition}[Central fiber and quotient defect]
The central additive fiber is
\[
 \mathcal F_{10}
 =\{(5+t,5-t):t\in\mathbb Z,\ |t|\leq4\}.
\]
In the chart \(n=\rho_9(n)+9q_9(n)\), with
\(\rho_9(n)\in\{1,\ldots,9\}\), the quotient defect relative
to the center is defined as
\[
 v(t)=q_9(25)-q_9(25-t^2),
\]
when both products have the same nonadic residue.
\end{definition}

\begin{lemma}[Minimal central vacancy]
The only points of \(\mathcal F_{10}\) that preserve the
multiplicative residue of the center are
\[
 (5,5),\qquad (8,2),\qquad (2,8).
\]
At the two noncentral points, the quotient defect is exactly one.
\end{lemma}
\begin{proof}
The sum is constantly ten, whereas
\[
 (5+t)(5-t)=25-t^2.
\]
Preservation of the residue is equivalent to \(t^2\equiv0\pmod9\), that is,
to \(3\mid t\). In the integer interval \([-4,4]\), the solutions are
\(t=-3,0,3\). For the two nonzero solutions,
\[
 25=7+9\cdot2,\qquad16=7+9\cdot1.
\]
The residue seven is preserved and one unit of quotient is lost. There is
no other admissible noncentral displacement in that domain, so the
positive defect is minimal. The interchange \(t\mapsto-t\) permutes the
two outputs and preserves the defect.
\end{proof}

Vacancy here is not a numerical zero: it is a difference between two
states with the same residual reading and different quotients. If only
the residue were preserved, that difference would disappear.
Orientation also distinguishes the two output points. Neither the quotient
defect nor that orientation is automatically identified with an electric
charge or a mass; their physical realization requires the corresponding
readers.

\subsection{Commutator norm of the arithmetic projections}

The electronic prefactor is obtained without using \(\alpha\), torsion,
or a dimensional unit. The domain is the third realization of APP,
\(\mathcal A_3=\{1,\ldots,9\}^3\), equipped with the uniform measure.
We define
\[
 \Sigma(x)=\rho_9(x_1+x_2+x_3),\qquad
 \Pi(x)=\rho_9(x_1x_2x_3).
\]
The conditional expectations \(P_\Sigma\) and \(P_\Pi\) are the
orthogonal projections onto functions constant on each additive
and multiplicative fiber, respectively. Thus, for a function \(f\),
\[
 (P_\Sigma f)(x)
 =\frac{1}{|\Sigma^{-1}(\Sigma(x))|}
   \sum_{y:\,\Sigma(y)=\Sigma(x)}f(y),
\]
and analogously for \(P_\Pi\). The two projections act on the
same space; they do not represent two independent arithmetic supports.

\begin{lemma}[Two projections and principal angles]
If \(s\in[0,1]\) is a singular value of the pairing between
orthonormal bases of the ranges of two projections \(P,Q\), its nontrivial
block contributes \(s\sqrt{1-s^2}\) to
\(\|[P,Q]\|\).
\end{lemma}
\begin{proof}
On the corresponding principal plane, one can write
\[
 P=\begin{pmatrix}1&0\\0&0\end{pmatrix},\qquad
 Q=\begin{pmatrix}
 s^2&s\sqrt{1-s^2}\\
 s\sqrt{1-s^2}&1-s^2
 \end{pmatrix}.
\]
Their commutator is
\[
 [P,Q]=s\sqrt{1-s^2}
 \begin{pmatrix}0&1\\-1&0\end{pmatrix}.
\]
The final matrix is orthogonal, so its norm is one. Common
or orthogonal blocks have zero commutator. The orthogonal decomposition
into principal planes gives the claim and reduces the total norm to the maximum
of these contributions.
\end{proof}

\begin{proposition}[Exact incompatibility norm]
On \(\ell^2(\mathcal A_3)\),
\begin{equation}
 \|[P_\Sigma,P_\Pi]\|=\frac{\sqrt3}{4}.
 \label{eq:el-prefactor}
\end{equation}
\end{proposition}
\begin{proof}
Let \(N_{ab}=|\Sigma^{-1}(a)\cap\Pi^{-1}(b)|\). Integer enumeration
of the three symbols yields
\[
 N=9\begin{pmatrix}
 0&1&1&0&1&1&0&1&4\\
 1&0&1&1&0&1&1&0&4\\
 1&0&2&0&1&2&0&0&3\\
 0&1&1&0&1&1&0&1&4\\
 1&0&1&1&0&1&1&0&4\\
 0&0&2&1&0&2&0&1&3\\
 0&1&1&0&1&1&0&1&4\\
 1&0&1&1&0&1&1&0&4\\
 0&1&2&0&0&2&1&0&3
 \end{pmatrix}.
\]
Each row sums to \(81\); the column sums are
\[
 (m_1,\ldots,m_9)=(36,36,108,36,36,108,36,36,297).
\]
The normalized indicator functions of the fibers have pairing
matrix \(C_{ab}=N_{ab}/\sqrt{81m_b}\). Consequently,
\[
 M=CC^{\mathsf T},\qquad
 M_{ac}=\sum_{b=1}^9\frac{N_{ab}N_{cb}}{81m_b}
\]
is a rational matrix whose spectrum consists of the squares of the
required singular values.

Exact multiplication shows that \((1,\ldots,1)\),
\((1,-1,0)^3\), and \((1,1,-2)^3\) are eigenvectors with eigenvalues
\(1\), \(1/4\), and \(7/132\), respectively. The subspace of
vectors supported on \(\{3,6,9\}\) whose sum is zero has dimension
two and eigenvalue \(1/18\). The zero-sum vectors supported on
\(\{1,4,7\}\) or \(\{2,5,8\}\) form a four-dimensional
kernel. These subspaces are independent and their dimensions sum to
nine. The complete spectrum is therefore established:
\[
 \operatorname{spec}(M)=
 \left\{1,\frac14,\frac7{132},\frac1{18},\frac1{18},0,0,0,0\right\}.
\]
The maximum of \(\lambda(1-\lambda)\) on this set is \(3/16\),
attained at \(\lambda=1/4\). The preceding lemma gives
\(\sqrt{3/16}=\sqrt3/4\).
\end{proof}

The mode \((1,-1,0)^3\) is precisely the distribution of values of the
tritic phase selector in the nonadic chart. The scalar norm retains
a measure of incompatibility, but does not by itself preserve the
orientation of that mode or reconstruct the two projectors. That memory
remains in the route state accompanying the scalar reading.

\begin{proposition}[Algebra of the local electronic block]
On the principal plane attaining \eqref{eq:el-prefactor}, let
\[
 P=\begin{pmatrix}1&0\\0&0\end{pmatrix},\qquad
 Q=\begin{pmatrix}1/4&\sqrt3/4\\\sqrt3/4&3/4\end{pmatrix},\qquad
 S=2P-I,\qquad J=\frac4{\sqrt3}[P,Q].
\]
Then \(S^2=I\), \(J^2=-I\), \(SJ=-JS\), and the real algebra
generated by \(S,J\) is \(M_2(\mathbb R)\), a realization of
\(\mathrm{Cl}_{1,1}\). Moreover,
\[
 n_\pm=\frac{S\pm J}{2},\qquad
 n_\pm^2=0,\qquad n_+n_-+n_-n_+=I.
\]
\end{proposition}
\begin{proof}
We have \(S=\operatorname{diag}(1,-1)\) and
\(J=\bigl(\begin{smallmatrix}0&1\\-1&0\end{smallmatrix}\bigr)\).
Multiplication verifies the three relations. The matrices
\(I,S,J,SJ\) are linearly independent and form a basis of
\(M_2(\mathbb R)\), proving the claim about the algebra.
Finally, \((S\pm J)^2=S^2+J^2\pm(SJ+JS)=0\), and the sum of the
cross products is \((S^2-J^2)/2=I\).
\end{proof}

In this block, an elliptic direction, a hyperbolic direction, and two
parabolic null directions coexist. The result describes a local
matrix representation of the arithmetic of the two sheets. It does not identify all of APP
with \(M_2(\mathbb R)\), or with a quantum Latin square; such
claims would require other domains and other relations. Nor does the
exponential belong exclusively to the parabolic regime: it can
be applied to any of the generators of this block.

\subsection{Transfer between sheets and coefficients of the basal stage}

The primitive tritic calendar visits the modes
\((\Sigma,\Pi,\Pi)\). Counting consecutive pairs, including the
closure of the block, yields one transition \(\Sigma\to\Pi\), one
\(\Pi\to\Pi\), and one \(\Pi\to\Sigma\). Their primitive incidence is
\[
 F_{\mathrm{av}}=\begin{pmatrix}0&1\\1&1\end{pmatrix}.
\]
This matrix realization is obtained after the calendar. It is not
used to replace the coinductive reader that has already generated \(\varphi\).
Its characteristic polynomial is \(X^2-X-1\), and recognition of the
two roots as \(\varphi\) and \(-\varphi^{-1}\) describes the action of
that output in modal transfer. The quadratic action has eigenvalues
\(\varphi^2,-1,\varphi^{-2}\); the last is the positive stable mode.
The regional mark \(\pi\), applied once, gives
\(\pi/\varphi^2\).

\begin{lemma}[Interchange of the two orientations]
Let \(\mathcal R(x,y)=x/y^2\) and \(J(x,y)=(y,x)\), for \(x,y>0\).
Then
\begin{equation}
 J^2=I,\qquad
 \mathcal R(\pi,\varphi)=\frac{\pi}{\varphi^2},\qquad
 (\mathcal R\circ J)(\pi,\varphi)=\frac{\varphi}{\pi^2}
 =\frac1{\varphi^3(\pi/\varphi^2)^2}.
 \label{eq:el-bisagra}
\end{equation}
\end{lemma}
\begin{proof}
Applying the interchange twice gives the identity. Moreover,
\(\varphi^3(\pi/\varphi^2)^2=\pi^2/\varphi\); taking its inverse proves
the last equality.
\end{proof}

The electronic reader uses the orientation \(\varphi/\pi^2\),
followed by the exponential character. The preceding identity explains the
interchange, but does not replace the rule that chooses that character or
by itself produce the integers of the correction. These have a different provenance.

The phase-compatible central return has minimal length
twenty-seven. With the phase and sheet record, its finite space is
\[
 \mathcal K_e=\mathbb Z/9\mathbb Z\times\mathbb F_3.
\]
The five regions of \(\pi\) are distinct genealogical identities,
not five interchangeable copies of a visible expansion. In the regional
gauge fixed by the construction, their injection into the return occupies
the positions
\begin{equation}
 (4,0),\quad(5,0),\quad(6,0),\quad(6,1),\quad(6,2).
 \label{eq:el-cinco-lugares}
\end{equation}
The transverse orientation, mutated return, and three positive
regions remain separated by that record. Admissible changes
of gauge transport the labels jointly; they do not reselect
the positions from an electronic value.

\begin{proposition}[Deficit and retained memory]
Given the central return and the preceding regional injection, the oriented
boundary deficit and torsional memory are
\begin{equation}
 n_e=-|\mathcal K_e\setminus\iota(\mathscr R_\pi)|=-22,
 \qquad
 m_e^{\mathrm{tor}}=|\mathscr R_\pi\times\mathbb F_3|=15.
 \label{eq:el-selector-enteros}
\end{equation}
\end{proposition}
\begin{proof}
The five positions in \eqref{eq:el-cinco-lugares} are distinct.
Thus, \(|\mathcal K_e|=9\cdot3=27\), and the complement of the image
has cardinality \(27-5=22\). The orientation convention records
the deficit removed from the boundary as negative. The retained memory has
three sheets per region, so its positive cardinality is
\(5\cdot3=15\).
\end{proof}

The same pair has the integer check supplied by the collection's APP chart:
\begin{equation}
 396=18+72+108+198,
 \qquad \frac{396}{18}=22,
 \qquad \frac{270}{18}=15=3\cdot5=\binom62.
 \label{eq:el-censo}
\end{equation}
The weight \(18\) belongs to the central \(2\times2\) kernel, and \(270\)
to the declared global channel. These equalities check the normalization
of the count; the return argument also fixes the signs. The
quotient of the alternating rings,
\(396/(18+108)=22/7\), is rational and is not identified with \(\pi\).
Nor does \(\binom62\) mean \(6/2\): it counts the unordered pairs
of a six-element set.

Evaluation of the deficit in the third symmetric realization uses
\(\alpha^3\). This power is the product of three copies of the
closure scalar already obtained; it does not assert a universal rule
associating any dimension with the same power of \(\alpha\).
The specified domain and reader are essential to the composition.

\subsection{Global torsion and the basal coordinate}

\begin{definition}[Electronic torsional normalization]
The reduced defect \(d_4\) and the complete reader
\(\Delta_4\) are distinguished:
\begin{equation}
 d_4=\frac{\pi-e\log\pi}{270},\qquad
 \Delta_4=d_4\left(1+\frac{\pi}{729}\right).
 \label{eq:el-delta}
\end{equation}
The denominator \(270\) corresponds to the census channel, whereas the
factor \(1+\pi/729\) retains the coupling with the cubic cell of
\(729\) positions. The electronic reading uses \(\Delta_4\) both in
its basal stage and in the return exponent.
\end{definition}

This definition acts on previously generated \(\pi\) and \(e\). The
logarithm is a subsequent operation in their real realization, not a
retrospective input of conventional values. Torsion is global:
it is not determined independently for each section by an observed
mass.

\begin{lemma}[The two normalizations are not interchangeable]
We have
\begin{equation}
 \Delta_4-d_4=\frac{\pi}{729}d_4.
 \label{eq:el-delta-diferencia}
\end{equation}
In particular, \(d_4>0\) and \(\Delta_4>d_4\) for the coordinates
under consideration.
\end{lemma}
\begin{proof}
The identity follows by distributivity. For positivity, the function
\(f(x)=x-e\log x\), on \(x>0\), has derivative \(1-e/x\) and
second derivative \(e/x^2>0\). Its unique minimum is at \(x=e\), where
it is zero. Since \(\pi\ne e\), we have \(f(\pi)>0\). Both factors
in \eqref{eq:el-delta} are positive and \(1+\pi/729>1\).
\end{proof}

\begin{remark}[Preservation of the normalization]
Removing the second factor in \eqref{eq:el-delta} changes the reader; it
is not an algebraic simplification. If \(d_4\) is introduced as an
auxiliary coordinate, the same object must be written as
\(\Delta_4=(1+\pi/729)d_4\) in every occurrence. Keeping the
integers \(15\) and \(6\) and replacing only \(\Delta_4\) by \(d_4\)
produces a different formula.
\end{remark}

\begin{definition}[Basal electronic coordinate]
The composition of incompatibility, modal transfer, deficit, and retained memory
is
\begin{equation}
 \varepsilon_e^{(0)}
 =\frac{\sqrt3}{4}
  \left[\exp\!\left(\frac{\varphi}{\pi^2}\right)-22\alpha^3\right]
  (1+15\Delta_4).
 \label{eq:el-basal}
\end{equation}
\end{definition}

The equation combines three operations of different kinds. The norm
\(\sqrt3/4\) quantifies the incompatibility of the sheet
projections; its origin is a principal angle of the APP fibers.
The exponential acts on the orientation \(\varphi/\pi^2\) of the
transfer reader. The cubic subtraction records the \(22\) complementary
positions of the marked return, and the torsional factor preserves the
\(15\) components corresponding to three sheets per region. The
composition preserves these domains even though its final evaluation is scalar.

The presence of \(\pi\), \(\varphi\), and \(e\) thus has operational
meanings. Euler's number enters through the exponential and
the logarithmic normalization of \(\Delta_4\); \(\varphi\) and \(\pi\)
enter through oriented transfer. The constant \(\alpha\)
acts after its own construction through the cubic closure
term. The chronological return must still act on this basal
coordinate. Separating the stages makes it possible to explain what changes when
an orientation, normalization, or record is modified, while preserving the central
structure that identifies the electronic section.

The parentheses in this equation are part of its content. The term
\(-22\alpha^3\) is subtracted \emph{outside} the exponential; the memory
\(1+15\Delta_4\) multiplies the entire bracket. Placing the deficit in
the exponent or applying a different torsion to it would alter the composition.

\begin{proposition}[Basal positivity]
For \(0<\alpha<10^{-2}\), \(\varphi>0\), and the normalization
\eqref{eq:el-delta}, we have \(\varepsilon_e^{(0)}>0\).
\end{proposition}
\begin{proof}
The exponential is greater than one, and
\(22\alpha^3<22\cdot10^{-6}<1\). The bracket is therefore positive.
Also, \(1+15\Delta_4>1\) and \(\sqrt3/4>0\).
\end{proof}

The subsequent evaluation of \eqref{eq:el-basal} begins with
\[
 \varepsilon_e^{(0)}=0.5109989224880660725\ldots.
\]
This number is still a dimensionless coordinate and is not the
final stage: the return memory of the route itself remains to be applied.

\subsection{Two readers of the electronic return}

Electronic memory is represented by a tritic word of length
\(108\) and a signed twelve-phase word. To make the exponent
calculation reproducible, we specify the chronological projection
used. Initially, let
\((r_0,c_0)=(5,5)\) and \((h_0,v_0)=(1,1)\). At each step, the selector
repeats \((+1,-1,0)\). With orientation fixed within each
twenty-seven-step block, the visible rule is
\[
 \begin{array}{c|c|c}
 \text{mode}&\text{integer read}&\text{cellular transport}\\\hline
 +1&r+c&c\mapsto1+(c-1+h)\bmod9\\
 -1&rc&r\mapsto1+(r-1+v)\bmod9\\
 0&9&\text{no displacement}
 \end{array}
\]
The emitted symbol is \(\rho_9(\text{integer read})\bmod3\).
After steps \(27,54,81\), \(h\) and \(v\) are simultaneously
reversed. This rule computes the observable projection; the
state update additionally preserves residue, quotient,
carry, sheet, boundary, and memory. The record of
\(108\) symbols is not identified with the complete state.

In each complete three-step block, each coordinate advances
once; nine such blocks are needed to return to the center in
the same phase. The minimal length of this recorded return is
therefore \(27\). If phase is forgotten, the cell is already reached after
step \(26\), before the neutral threshold completing the block; that is not
the typed return counted by \(\mathcal K_e\). The explicit reading
of the first two return blocks gives
\begin{equation}
 a=100000220,\qquad b=120200000,\qquad
 \gamma_e=(a^3b^3)^2.
 \label{eq:el-palabra}
\end{equation}
Here, powers denote word concatenation. In particular,
\(\gamma_{e,t+54}=\gamma_{e,t}\). In the record of nine steps per
phase, the central orientations give
\begin{equation}
 \tau_e=(1,1,1,-1,-1,-1,1,1,1,-1,-1,-1).
 \label{eq:el-tau}
\end{equation}
The periodicity of these representations does not amount to restarting the
history that produces them.

\begin{definition}[Cyclic discrepancy reader]
For a word \(\gamma\in\mathbb F_3^{108}\), define
\[
 D_k(\gamma)=\#\{t\in\mathbb Z/108\mathbb Z:
                   \gamma_t\ne\gamma_{t+k}\}.
\]
On the domain of words whose period divides \(54\), the integer
reader is
\begin{equation}
 K_D(\gamma)=\left(D_{27}+D_{36},D_1,\frac{D_4-D_3}{2}\right).
 \label{eq:el-kd}
\end{equation}
\end{definition}

\begin{lemma}[Integrality and central evaluation]
The \(D_k\) on the preceding domain are even. For
\eqref{eq:el-palabra},
\[
 \begin{array}{c|rrrrr}
 k&1&3&4&27&36\\\hline
 D_k&54&56&68&48&32
 \end{array}
\]
and \(K_D(\gamma_e)=(80,54,6)\).
\end{lemma}
\begin{proof}
The map \(t\mapsto t+54\) pairs discrepancy indices without fixed
points and preserves their condition; each cardinality is even. For the
explicit word \(a^3b^3\) of length \(54\), the discrepancy
count for each displacement is, in the same order,
\(27,28,34,24,16\). Doubling the word doubles these five cardinalities.
Finally,
\[
 (D_{27}+D_{36},D_1,(D_4-D_3)/2)
 =(48+32,54,(68-56)/2)=(80,54,6).
\]
\end{proof}

The displacements \(27\) and \(36\) are the chronological
lifts of the \(90\)- and \(120\)-degree channels in a chart
where nine steps correspond to thirty degrees. The angular reading
explains that representation, but the integers are computed on the
word, without introducing a mass or a target real value.

\begin{definition}[Dodecaphasic window reader]
For \(\tau\in\{-1,0,1\}^{12}\), with cyclic indices, let
\begin{align*}
 C_k(\tau)&=\sum_{r=0}^{11}\tau_r\tau_{r+k},\\
 A_\ell(\tau)&=\sum_{r=0}^{11}
           \left(\sum_{j=0}^{\ell-1}\tau_{r+j}\right)^2,\\
 P_6(\tau)&=\sum_{r=0}^{5}\tau_r\tau_{r+6}.
\end{align*}
Define
\begin{equation}
 \Omega(\tau)=4A_3(\tau)-3A_4(\tau),\qquad
 K_\Omega(\tau)=(\Omega(\tau),54,P_6(\tau)).
 \label{eq:el-komega}
\end{equation}
The \(54\) in this reader records the declared chronological half-turn.
It is not confused with the local APP jump of value four.
\end{definition}

\begin{lemma}[Primitive contrast and correlation form]
Among integer linear contrasts \(uA_3+vA_4\) that vanish
whenever \(A_3=3a\) and \(A_4=4a\), the primitive pair, with
orientation \(90-120\), is \((u,v)=(4,-3)\). Moreover,
\begin{equation}
 \Omega=-2C_1-4C_2-6C_3.
 \label{eq:el-omega-corr}
\end{equation}
\end{lemma}
\begin{proof}
Vanishing for every \(a\) requires \(3u+4v=0\). By coprimality,
the solutions are \((u,v)=n(4,-3)\), and orientation fixes the sign
of the primitive generator. Expanding the squares gives
\[
 A_\ell=\ell C_0+2\sum_{k=1}^{\ell-1}(\ell-k)C_k.
\]
Therefore,
\(A_3=3C_0+4C_1+2C_2\) and
\(A_4=4C_0+6C_1+4C_2+2C_3\). Their combination \(4A_3-3A_4\)
cancels \(C_0\) and yields \eqref{eq:el-omega-corr}.
\end{proof}

The uniqueness proved belongs to that class of integer linear
contrasts; it does not by itself exclude every nonlinear functional of the windows.
Specifying the functional class avoids turning an identity
of two channels into a claim of universal selection.

\begin{theorem}[Return triple of the electronic section]
The two readers applied to the records of \(\Gamma_e\)
coincide and give
\begin{equation}
 K_D(\gamma_e)=K_\Omega(\tau_e)=(80,54,6).
 \label{eq:el-triple}
\end{equation}
Their evaluation in the already generated coordinates is
\begin{equation}
 \kappa_e=80-54\alpha+6\Delta_4.
 \label{eq:el-kappa}
\end{equation}
\end{theorem}
\begin{proof}
The first reader was computed above. For \(\tau_e\), the antipodal
product is always one, and hence \(P_6=6\). The cyclic count gives
\(C_1=4\), \(C_2=-4\), \(C_3=-12\); thus
\[
 \Omega=-2\cdot4-4(-4)-6(-12)=80.
\]
Equivalently, \(A_3=44\), \(A_4=32\), and
\(4\cdot44-3\cdot32=80\). Both triples are therefore
\((80,54,6)\). Evaluating the functional
\((A,B,C)\mapsto A-\alpha B+\Delta_4C\) proves the last formula.
\end{proof}

This coincidence concerns the central section. It does not identify
the two readers on every route domain: one compares tritic
symbols and the other compares windows of a signed record. Nor is
\(80\) obtained from the period of a word assigned to Euler's
number. Its origin in this composition is the triple of the electronic
route.

\subsection{Complete electronic composition and normalization}

\begin{definition}[Multiplicative return]
The action of return memory on the basal coordinate is
\begin{equation}
 \varepsilon_e^\beta
 =\varepsilon_e^{(0)}\exp\!\left(\kappa_e\log R_{\mathrm{act}}\right).
 \label{eq:el-beta}
\end{equation}
\end{definition}

\begin{theorem}[Positive electronic factorization]
With the declared readers, orientation, and normalization,
\eqref{eq:el-beta} coincides with \eqref{eq:el-formula-completa}. All
its factors are dimensionless and the resulting coordinate is positive.
\end{theorem}
\begin{proof}
The incompatibility norm is \eqref{eq:el-prefactor}; modal transfer,
deficit, and retained memory compose \eqref{eq:el-basal}. The
triple \eqref{eq:el-triple} produces \eqref{eq:el-kappa}, and the ratio of
the two action sections is \eqref{eq:el-ratio}. Substituting those
three identities into \eqref{eq:el-beta} yields exactly
\eqref{eq:el-formula-completa}. The basal stage is positive, as is the
exponential of a real number. The common factor with the dimension of
action has already canceled in forming \(R_{\mathrm{act}}\).
\end{proof}

The evaluation of this composition begins with
\begin{equation}
 \varepsilon_e^\beta
 =0.5109989506900008086082571648\ldots.
 \label{eq:el-beta-decimal}
\end{equation}
The digits are a subsequent evaluation of the formula, not data used
to select \(22\), \(15\), or \((80,54,6)\). The numerical
representation of the basal coordinate by a finite prefix introduces only its
truncation error; it does not modify the exact formula defining it.

The effect of a different torsional normalization can be located
without comparison to a measurement. Write
\(B=\frac{\sqrt3}{4}[\exp(\varphi/\pi^2)-22\alpha^3]\),
\(R=R_{\mathrm{act}}\), and
\[
 \mathcal E(d)=B(1+15d)R^{80-54\alpha+6d}.
\]
For the two defects in \eqref{eq:el-delta},
\begin{equation}
 \frac{\mathcal E(\Delta_4)}{\mathcal E(d_4)}
 =\frac{1+15\Delta_4}{1+15d_4}
   R^{6(\Delta_4-d_4)}>1.
 \label{eq:el-normalizacion-beta}
\end{equation}
The first torsional inequality and \(R>1\) prove that they are not the same
coordinate. If only the defect in the exponent is changed, keeping the
complete basal coordinate, the ratio of the two outputs is
\(R^{6(\Delta_4-d_4)}\), again different from one. A very small
difference is still a difference of formula; it cannot be attributed to
a mere conversion of units.

\subsection{Dimensional realization: energy, mass, and chronological scale}

The electronic coordinate finally acts on a dimensional line.
Let \(\mathcal L_E\) be an energy line with positive basis
\(\mathcal U_E\). The realization is
\begin{equation}
 E_e^{(0)}=\varepsilon_e^{(0)}\mathcal U_E,\qquad
 E_e^\beta=\varepsilon_e^\beta\mathcal U_E,
 \qquad m_e^\beta=\frac{E_e^\beta}{c_{\mathrm{int}}^2}.
 \label{eq:el-dimensional}
\end{equation}
The last expression belongs to the mass line when
\(c_{\mathrm{int}}\) has the dimension of velocity in the declared
realization. In the comparative chart \(\mathcal U_E=1\,\mathrm{MeV}\),
\eqref{eq:el-beta-decimal} expresses an energy in MeV; the corresponding
mass is expressed in \(\mathrm{MeV}/c^2\). It is common
to abbreviate both quantities by taking \(c=1\), but that abbreviation must
not conceal the dimensional map.

The action scale also admits an action--clock realization. If
\(t_0\) is a positive temporal section, the \(108\)-step clock
defines
\begin{equation}
 \omega_0=\frac{2\pi}{108t_0},\qquad
 E_{\mathrm{clk}}=\hbar_{\mathrm{ret}}\omega_0,\qquad
 m_{\mathrm{clk}}=E_{\mathrm{clk}}c_{\mathrm{int}}^{-2}.
 \label{eq:el-cuanto-reloj}
\end{equation}
Action times frequency has the dimension of energy. This chronological
quantum is a common dimensional section; it is not, by definition,
the electronic section. The route word and its factors are not
obtained from the uniformity of the clock.

\begin{proposition}[Cancellation of the common scale]
The factor \(10^{-34}\mathcal U_S\) in
\eqref{eq:el-secciones-accion} is not part of \(\kappa_e\). On
the same energy line,
\begin{equation}
 \frac{E_e^\beta}{E_e^{(0)}}=R_{\mathrm{act}}^{\kappa_e}.
 \label{eq:el-razon-energia}
\end{equation}
The result is independent of the choice of the common dimensional
basis.
\end{proposition}
\begin{proof}
In forming \(\hbar_{\mathrm{pre}}/\hbar_{\mathrm{ret}}\), the factor
\(10^{-34}\mathcal U_S\) cancels exactly. Dividing the two
energy expressions in \eqref{eq:el-dimensional} also cancels
\(\mathcal U_E\), and \eqref{eq:el-beta} gives the stated ratio.
\end{proof}

Thus, a coordinate already expressed in MeV is not multiplied again
by \(10^{-34}\). If \(E_{\mathrm{clk}}\) is chosen as the basis of
\(\mathcal L_E\), the coordinates must be transported by the
change-of-basis factor; retaining the number from the one-MeV chart
unchanged would identify two bases without proof. This is the precise point
at which the action scale enters the absolute realization and ceases to
enter the relative correction.

There is also a linear realization that makes the separation
between clock and genealogy visible. On
\(\ell^2(\mathbb Z/108\mathbb Z;\mathbb R)\), the cyclic
shift \(S\) has fixed space
\[
 \ker(S-I)=\mathbb R n_0,\qquad
 n_0=\frac1{\sqrt{108}}\sum_{t=0}^{107}\delta_t.
\]
If \(g_{\Gamma_e}\) is the route symbol in the free linearization
of genealogies and \(\mathcal U_M\) is a basis of the mass line,
the rank-one map is defined by
\begin{equation}
 \mathfrak C_e:
 \mathbb R n_0\otimes\mathbb R g_{\Gamma_e}\longrightarrow\mathcal L_M,
 \qquad
 \mathfrak C_e(n_0\otimes g_{\Gamma_e})
 =\varepsilon_e^\beta\mathcal U_M.
 \label{eq:el-rango-uno}
\end{equation}
The map assembles the preceding factorization once its bases have been declared.
The uniform vector \(n_0\) does not select \(\varepsilon_e^\beta\):
that coordinate comes from APP, the oriented return, and its readers.
An algebraic realization of the null mode is thus distinguished from an
identification of the electron with the clock quantum.

The internal conclusion is the positive factorization of a central
section, with its operations, memory, and normalizations made explicit. Assignment
to a physical electron adds the charge representation,
spin, coupling, and corresponding unit chart; it does not
follow exclusively from a decimal coincidence. Metrological
comparison is presented after the exceptional prolongation and does not alter
the inputs of any of the constructions performed here.

\subsubsection{Structural dependence of the electronic coordinate}

The electronic equation expresses a hierarchical composition. Its first
factor, \(\sqrt3/4\), belongs to the central matrix structure and fixes
a geometric normalization. The exponent \(\varphi/\pi^2\) uses the
already generated regional coordinates in the indicated orientation of the
areal--volumetric reader. The term \(22\alpha^3\) incorporates the dodecaphasic
coupling into the cubic correction. The factor \(1+15\Delta_4\) preserves
the global torsional normalization. Finally, the quotient of the action
sections acts with the exponent determined by the return readers.
Each operation thus has a different antecedent within the same
construction.

The articulation becomes more precise when its dependencies are considered.
The central norm can be calculated before the regional coordinates
are evaluated. The exponential and cubic correction require those
coordinates and the closure of \(\alpha\). Multiplicative memory additionally
requires the action sections and the \(108\)-position trajectory.
Knowing the basal expression therefore does not amount to having evaluated
the complete coordinate. The succession of operations does not add an
external mass to the point \((5,5)\); it develops the readers assigned to its
persistent section.

The distinction between scale and structure also provides an
interpretation of the physical comparison. Mass is the dimensional
realization of the complete coordinate, while the centre, involution,
and vacancy describe the relations that produce it. A common change
of unit modifies the dimensional coordinates, but leaves the dimensionless
ratios and the proved transport identities unchanged.
This is the sense in which the electron preserves its structural character
when expressed relative to a chronological scale.



\par
\subsection{Lorentzian transport and electromagnetic representation}

The electronic section provides a structure and a coordinate
before its application to an interaction. The Lorentzian realization
of cubic incidence determines, under the stated temporal convention,
the space of forms in which that interaction can be expressed.
To formulate it, an oriented local realization \((\mathcal M,g)\)
of dimension four is fixed, with tangent fibres identified with the model
\(Q_\square\). A connection \(\nabla\) with
\(\nabla g=0\), a potential one-form \(A\in\Omega^1(\mathcal M)\),
its curvature \(F=dA\), and a charge representation of the section
are also specified. These data are the arguments of the physical representation
that follows.

Hodge duality satisfies \(\star^2=-I\) on two-forms.
The electric--magnetic separation corresponds to a temporal choice
in the \(1+3\) decomposition; \(F\) and \(\star F\) preserve their
relationship within the same space of two-forms. Neutrality of a
charge representation means \(q=0\), whereas the section may
retain orientation, phase, and internal transport. The absence of
electric coupling in that representation has a different content
from the absence of phase information.

In this dimensional realization, let \(m_e>0\), let \(q\) be the charge,
let \(\gamma(s)\) be a parametrized trajectory, and let \(u=\dot\gamma\)
be its four-velocity. The Lorentz force is expressed by
\begin{equation}
 m_e\nabla_u u^\mu=qF^\mu{}_{\nu}u^\nu.
 \label{eq:lorentz-electron}
\end{equation}
The antisymmetry of \(F\) implies
\(F_{\mu\nu}u^\mu u^\nu=0\). Consequently,
\[
 \frac{d}{ds}g(u,u)=2g(\nabla_u u,u)=0.
\]
The normalization of the trajectory is preserved provided that \(\nabla\)
is metric and \eqref{eq:lorentz-electron} holds. The mass used
is the dimensional realization of the preceding electronic section;
the coupling parameter is identified in the adopted electromagnetic
convention. The concrete deduction of the potential and of a field
configuration constitutes an additional problem with its own
boundary conditions. The result presented here determines the domain, the
representation data, and the conservation implied by its equation.

TPK route inversion and the orientation of an electronic line
then offer a precise comparison: each has its involution,
transport, and representation. The work of Wheeler and Feynman
places the question historically; the equivalence of two realizations
is decided through their maps and compatibility conditions.
This arrangement keeps the electron's arithmetic origin connected
to the formulation of its interaction, making the function
of each stage explicit.

